
\documentclass{article}
\usepackage{colortbl}
\usepackage{makecell}
\usepackage{multirow}
\usepackage{supertabular}

\begin{document}

\newcounter{utterance}

\centering \large Interaction Transcript for game `cladder', experiment `full\_v1.5\_default', episode 9786 with score{-}Qwen3.5{-}9B{-}sp.
\vspace{24pt}

{ \footnotesize  \setcounter{utterance}{1}
\setlength{\tabcolsep}{0pt}
\begin{supertabular}{c@{$\;$}|p{.15\linewidth}@{}p{.15\linewidth}p{.15\linewidth}p{.15\linewidth}p{.15\linewidth}p{.15\linewidth}}
    \# & \multicolumn{2}{c}{Player} && \multicolumn{2}{c}{Game Master} \\
    \hline

    \theutterance \stepcounter{utterance}  
    & & & \multicolumn{4}{p{0.6\linewidth}}{
        \cellcolor[rgb]{0.9,0.9,0.9}{
            \makecell[{{p{\linewidth}}}]{
                \texttt{\tiny{[P1$\langle$GM]}}
                \texttt{You are an expert in causal inference. The following question is not a typical commonsense query, but rather a meticulously designed question created by a professor specializing in causal inference, intended to assess the students' mastery of the course content.} \\
\texttt{Imagine a self{-}contained, hypothetical world with only the following conditions, and without any unmentioned factors or causal relationships: Qwiu has a direct effect on jyka and kwox. Jyka has a direct effect on yupt. Yupt has a direct effect on kwox. Qwiu is unobserved. We know that qwiu causes jyka. jyka causes yupt. qwiu or yupt causes kwox. We observed an individual is not qwiu. Would an individual is not kwox if not jyka instead of jyka?} \\
\texttt{Start your answer with 'Yes' or 'No', followed by additional reasoning or evidence to support your explanation.} \\
            }
        }
    }
     \\ \\

    \theutterance \stepcounter{utterance}  
    & \multicolumn{4}{p{0.6\linewidth}}{
        \cellcolor[rgb]{0.561,0.737,0.561}{
            \makecell[{{p{\linewidth}}}]{
                \texttt{\tiny{[P1$\rangle$GM]}}
                \texttt{No} \\
\\ 
\texttt{**Reasoning:**} \\
\\ 
\texttt{To determine the answer, we must analyze the causal structure and the specific counterfactual query within the constraints of the hypothetical world.} \\
\\ 
\texttt{1.  **Causal Structure Identification:**} \\
\texttt{    *   **Nodes:** Qwiu (Unobserved), Jyka, Yupt, Kwox.} \\
\texttt{    *   **Edges (Direct Effects):**} \\
\texttt{        *   \$Qwiu \textbackslash{}rightarrow Jyka\$ (Explicitly stated: "Qwiu causes jyka").} \\
\texttt{        *   \$Jyka \textbackslash{}rightarrow Yupt\$ (Explicitly stated: "jyka causes yupt").} \\
\texttt{        *   \$Yupt \textbackslash{}rightarrow Kwox\$ (Explicitly stated: "Yupt has a direct effect on kwox").} \\
\texttt{        *   \$Qwiu \textbackslash{}rightarrow Kwox\$ (Implied by "Qwiu has a direct effect on... kwox").} \\
\texttt{        *   \$Yupt \textbackslash{}rightarrow Kwox\$ (Explicitly stated: "Yupt has a direct effect on kwox").} \\
\texttt{    *   **Ambiguity Resolution:** The prompt states "qwiu or yupt causes kwox." Given the explicit statement that "Yupt has a direct effect on kwox," and the general rule that unmentioned factors do not exist, we treat the relationship as \$Yupt \textbackslash{}rightarrow Kwox\$. The phrase "qwiu or yupt" likely describes} \\
            }
        }
    }
    & & \\ \\

    \theutterance \stepcounter{utterance}  
    & & & \multicolumn{2}{p{0.3\linewidth}}{
        \cellcolor[rgb]{0.9,0.9,0.9}{
            \makecell[{{p{\linewidth}}}]{
                \texttt{\tiny{[GM$|$GM]}}
                \texttt{No} \\
            }
        }
    }
    & & \\ \\

    \theutterance \stepcounter{utterance}  
    & & & \multicolumn{2}{p{0.3\linewidth}}{
        \cellcolor[rgb]{0.9,0.9,0.9}{
            \makecell[{{p{\linewidth}}}]{
                \texttt{\tiny{[GM$|$GM]}}
                \texttt{yes} \\
            }
        }
    }
    & & \\ \\

    \theutterance \stepcounter{utterance}  
    & & & \multicolumn{2}{p{0.3\linewidth}}{
        \cellcolor[rgb]{0.9,0.9,0.9}{
            \makecell[{{p{\linewidth}}}]{
                \texttt{\tiny{[GM$|$GM]}}
                \texttt{game\_result = LOSE} \\
            }
        }
    }
    & & \\ \\

\end{supertabular}
}

\end{document}
